\renewcommand{\tema}{Preguntas Post-Clase - Trabajo y Leyes de la Conservación}

\twocolumn
\section{PREGUNTAS POST-CLASE}

% -------------------------------------------------------
% BLOQUE 1 — Aplicación directa (P1-P5)
% -------------------------------------------------------

\begin{preguntabox}
\subsection{Pregunta 1}
Un técnico hala una turbina de aerogenerador a lo largo del piso con una fuerza de 500 N formando un ángulo de 60$^\circ$ con la horizontal. Si la desplaza 4 m, ¿cuánto trabajo realizó?

\begin{enumerate}
    \item[A)] 0 J
    \item[B)] 1\,000 J
    \item[C)] 1\,730 J
    \item[D)] 2\,000 J
\end{enumerate}
\end{preguntabox}

\begin{preguntabox}
\subsection{Pregunta 2}
Un ciclista profesional carga su bicicleta (10 kg) mientras camina 500 m en línea recta sobre el piso del velódromo. ¿Cuánto trabajo realizan sus brazos sobre la bicicleta?

\begin{enumerate}
    \item[A)] 5\,000 J
    \item[B)] 50\,000 J
    \item[C)] 0 J
    \item[D)] $-$5\,000 J
\end{enumerate}
\end{preguntabox}

\begin{preguntabox}
\subsection{Pregunta 3}
Una grúa de construcción desarrolla una potencia constante de 12 kW. ¿Cuánto trabajo total realiza en 5 minutos de operación continua?

\begin{enumerate}
    \item[A)] 2\,400 J
    \item[B)] 60\,000 J
    \item[C)] 3\,600\,000 J
    \item[D)] 300\,000 J
\end{enumerate}
\end{preguntabox}

\begin{preguntabox}
\subsection{Pregunta 4}
Un satélite de 8 kg tiene una energía cinética de 1\,600 J en órbita. ¿Con qué rapidez se desplaza?

\begin{enumerate}
    \item[A)] 10 m/s
    \item[B)] 20 m/s
    \item[C)] 200 m/s
    \item[D)] 400 m/s
\end{enumerate}
\end{preguntabox}

\begin{preguntabox}
\subsection{Pregunta 5}
El resorte del amortiguador de un camión de carga tiene $k = 50\,000$ N/m y se comprime 0.04 m al pasar sobre un tope. ¿Cuánta energía potencial elástica almacena?

\begin{enumerate}
    \item[A)] 2\,000 J
    \item[B)] 80 J
    \item[C)] 40 J
    \item[D)] 4 J
\end{enumerate}
\end{preguntabox}

% -------------------------------------------------------
% BLOQUE 2 — Combinación de conceptos (P6-P10)
% -------------------------------------------------------

\begin{preguntabox}
\subsection{Pregunta 6}
Un tractor arrastra un arado con una fuerza de 1\,000 N a 60$^\circ$ sobre la horizontal, moviéndose a velocidad constante de 4 m/s. ¿Qué potencia desarrolla el motor?

\begin{enumerate}
    \item[A)] 1\,000 W
    \item[B)] 2\,000 W
    \item[C)] 3\,460 W
    \item[D)] 4\,000 W
\end{enumerate}
\end{preguntabox}

\begin{preguntabox}
\subsection{Pregunta 7}
Un cañón de aire comprimido lanza un dardo de 2 kg a partir del reposo. Si le realiza 100 J de trabajo, ¿con qué velocidad sale el dardo?

\begin{enumerate}
    \item[A)] 5 m/s
    \item[B)] 10 m/s
    \item[C)] 50 m/s
    \item[D)] 100 m/s
\end{enumerate}
\end{preguntabox}

\begin{preguntabox}
\subsection{Pregunta 8}
Un clavadista de 70 kg salta desde una plataforma a 20 m de altura sobre la superficie del agua. Usando conservación de energía, ¿con qué rapidez entra al agua?

\begin{enumerate}
    \item[A)] 10 m/s
    \item[B)] 14 m/s
    \item[C)] 20 m/s
    \item[D)] 200 m/s
\end{enumerate}
\end{preguntabox}

\begin{preguntabox}
\subsection{Pregunta 9}
El resorte lanzador de un pinball tiene $k = 800$ N/m y se comprime 0.5 m. La bola tiene 2 kg. Suponiendo que toda la energía elástica se convierte en cinética, ¿con qué velocidad sale la bola?

\begin{enumerate}
    \item[A)] 5 m/s
    \item[B)] 10 m/s
    \item[C)] 20 m/s
    \item[D)] 100 m/s
\end{enumerate}
\end{preguntabox}

\begin{preguntabox}
\subsection{Pregunta 10}
Dos vagones de tren: el vagón A (10\,000 kg) se mueve a 4 m/s y choca con el vagón B (10\,000 kg) en reposo. La colisión es perfectamente inelástica. ¿Con qué velocidad se desplazan juntos?

\begin{enumerate}
    \item[A)] 4 m/s
    \item[B)] 2 m/s
    \item[C)] 1 m/s
    \item[D)] 0.5 m/s
\end{enumerate}
\end{preguntabox}

% -------------------------------------------------------
% BLOQUE 3 — Problemas desafiantes (P11-P15)
% -------------------------------------------------------

\begin{preguntabox}
\subsection{Pregunta 11}
Un snowboarder desciende una rampa nevada (sin fricción) desde una altura de 20 m. ¿Con qué rapidez llega al fondo?

\begin{enumerate}
    \item[A)] 10 m/s
    \item[B)] 20 m/s
    \item[C)] 40 m/s
    \item[D)] 200 m/s
\end{enumerate}
\end{preguntabox}

\begin{preguntabox}
\subsection{Pregunta 12}
El motor de una planta eólica consume 200 kJ para generar 130 kJ de energía eléctrica. ¿Cuál es su eficiencia?

\begin{enumerate}
    \item[A)] 35\%
    \item[B)] 54\%
    \item[C)] 65\%
    \item[D)] 154\%
\end{enumerate}
\end{preguntabox}

\begin{preguntabox}
\subsection{Pregunta 13}
Un tobogán acuático tiene 10 m de longitud y 8 m de altura. Un niño de 10 kg baja con una fricción constante de 30 N. ¿Con qué velocidad llega al fondo?

\begin{enumerate}
    \item[A)] 5 m/s
    \item[B)] 10 m/s
    \item[C)] 16 m/s
    \item[D)] 20 m/s
\end{enumerate}
\end{preguntabox}

\begin{preguntabox}
\subsection{Pregunta 14}
Un atleta paralímpico en silla de ruedas (masa total 90 kg) sube una rampa de 4 m de altura en 20 s. ¿Qué potencia promedio desarrolla?

\begin{enumerate}
    \item[A)] 18 W
    \item[B)] 72 W
    \item[C)] 180 W
    \item[D)] 3\,600 W
\end{enumerate}
\end{preguntabox}

\begin{preguntabox}
\subsection{Pregunta 15}
Un astronauta (100 kg) flota en la ISS sin moverse. Empuja un panel solar (20 kg) que sale a 3 m/s. ¿Con qué rapidez retrocede el astronauta?

\begin{enumerate}
    \item[A)] 3 m/s
    \item[B)] 1.5 m/s
    \item[C)] 0.6 m/s
    \item[D)] 15 m/s
\end{enumerate}
\end{preguntabox}

% -------------------------------------------------------
% RESPUESTAS
% -------------------------------------------------------

\newpage
\onecolumn
\section{RESPUESTAS Y RETROALIMENTACIÓN}

% --- P1 ---
\begin{respuestabox}
\subsection{Pregunta 1}
\textcolor{verdecorrecto}{\textbf{Respuesta correcta: B) 1\,000 J}}

\textbf{Justificación:}\\
\textbf{Fórmula utilizada:} Trabajo mecánico (Lección: Sección 1 --- Trabajo Mecánico)
$$W = F \cdot d \cdot \cos\theta = 500\,\text{N} \times 4\,\text{m} \times \cos 60^\circ = 500 \times 4 \times 0.5 = 1\,000\,\text{J}$$

\textbf{Respuestas incorrectas:}
\begin{itemize}
    \item \textbf{A) 0 J:} Confunde 60$^\circ$ con 90$^\circ$ --- solo habría trabajo nulo si la fuerza fuera perpendicular al desplazamiento.
    \item \textbf{C) 1\,730 J:} Usa $\cos 30^\circ \approx 0.866$ en lugar de $\cos 60^\circ = 0.5$ --- confunde el ángulo con su complementario.
    \item \textbf{D) 2\,000 J:} Calcula $F \times d$ sin aplicar $\cos\theta$ --- ignora el ángulo.
\end{itemize}
\end{respuestabox}

% --- P2 ---
\begin{respuestabox}
\subsection{Pregunta 2}
\textcolor{verdecorrecto}{\textbf{Respuesta correcta: C) 0 J}}

\textbf{Justificación:}\\
\textbf{Fórmula utilizada:} Trabajo mecánico (Lección: Sección 1 --- Trabajo Mecánico)
$$W = F \cdot d \cdot \cos\theta$$

Los brazos sostienen la bicicleta verticalmente, pero el desplazamiento es horizontal: $\theta = 90^\circ$ y $\cos 90^\circ = 0$, por lo que $W = 0$.

\textbf{Respuestas incorrectas:}
\begin{itemize}
    \item \textbf{A) 5\,000 J:} Usa $W = mgd$ con valores mezclados --- confunde desplazamiento horizontal con altura.
    \item \textbf{B) 50\,000 J:} Usa $W = mgd = 10 \times 10 \times 500$ --- ignora la dirección de la fuerza.
    \item \textbf{D) $-$5\,000 J:} Aplica signo negativo sin justificación; la fuerza no frena el movimiento.
\end{itemize}
\end{respuestabox}

% --- P3 ---
\begin{respuestabox}
\subsection{Pregunta 3}
\textcolor{verdecorrecto}{\textbf{Respuesta correcta: C) 3\,600\,000 J}}

\textbf{Justificación:}\\
\textbf{Fórmula utilizada:} Potencia (Lección: Sección 2 --- Potencia)
$$W = P \cdot t$$

\textbf{Conversiones:} $12\,\text{kW} = 12\,000\,\text{W}$ \quad y \quad $5\,\text{min} = 5 \times 60\,\text{s} = 300\,\text{s}$

\textbf{Sustitución:}
$$W = 12\,000\,\text{W} \times 300\,\text{s} = 3\,600\,000\,\text{J}$$

\textbf{Respuestas incorrectas:}
\begin{itemize}
    \item \textbf{A) 2\,400 J:} Divide $P/t$ en lugar de multiplicar.
    \item \textbf{B) 60\,000 J:} Multiplica $12\,000 \times 5$ (minutos) sin convertir a segundos.
    \item \textbf{D) 300\,000 J:} Ignora el prefijo kilo: usa $P = 1\,000$ W y calcula $1\,000 \times 300$.
\end{itemize}
\end{respuestabox}

% --- P4 ---
\begin{respuestabox}
\subsection{Pregunta 4}
\textcolor{verdecorrecto}{\textbf{Respuesta correcta: B) 20 m/s}}

\textbf{Justificación:}\\
\textbf{Fórmula utilizada:} Energía cinética (Lección: Sección 3 --- Energía Cinética)
$$E_c = \frac{1}{2}mv^2 \quad \Rightarrow \quad v = \sqrt{\frac{2E_c}{m}} = \sqrt{\frac{2 \times 1\,600\,\text{J}}{8\,\text{kg}}} = \sqrt{400\,\text{m}^2/\text{s}^2} = 20\,\text{m/s}$$

\textbf{Respuestas incorrectas:}
\begin{itemize}
    \item \textbf{A) 10 m/s:} Usa $v = \sqrt{E_c/m}$ sin el factor 2: $\sqrt{200} \approx 14$ m/s, redondeado.
    \item \textbf{C) 200 m/s:} Calcula $2E_c/m = 400$ pero olvida aplicar la raíz cuadrada.
    \item \textbf{D) 400 m/s:} Calcula $2E_c/m$ con error adicional y lo reporta directamente.
\end{itemize}
\end{respuestabox}

% --- P5 ---
\begin{respuestabox}
\subsection{Pregunta 5}
\textcolor{verdecorrecto}{\textbf{Respuesta correcta: C) 40 J}}

\textbf{Justificación:}\\
\textbf{Fórmula utilizada:} Energía potencial elástica (Lección: Sección 4 --- Energía Potencial)
$$E_e = \frac{1}{2}kx^2 = \frac{1}{2}(50\,000\,\text{N/m})(0.04\,\text{m})^2 = \frac{1}{2}(50\,000)(0.0016) = 40\,\text{J}$$

\textbf{Respuestas incorrectas:}
\begin{itemize}
    \item \textbf{A) 2\,000 J:} Calcula $kx = 50\,000 \times 0.04$ sin elevar al cuadrado ni aplicar $\frac{1}{2}$.
    \item \textbf{B) 80 J:} Calcula $kx^2 = 50\,000 \times 0.0016 = 80$ --- olvida el factor $\frac{1}{2}$.
    \item \textbf{D) 4 J:} Usa $x = 0.4$ m por error de decimal en la compresión.
\end{itemize}
\end{respuestabox}

% --- P6 ---
\begin{respuestabox}
\subsection{Pregunta 6}
\textcolor{verdecorrecto}{\textbf{Respuesta correcta: B) 2\,000 W}}

\textbf{Justificación:}\\
\textbf{Fórmula utilizada:} Potencia con fuerza a un ángulo (Lección: Secciones 1 y 2)
$$P = F\cos\theta \cdot v = 1\,000\,\text{N} \times \cos 60^\circ \times 4\,\text{m/s} = 1\,000 \times 0.5 \times 4 = 2\,000\,\text{W}$$

Solo la componente horizontal ($F\cos\theta$) realiza trabajo sobre el desplazamiento horizontal.

\textbf{Respuestas incorrectas:}
\begin{itemize}
    \item \textbf{A) 1\,000 W:} Calcula $F\cos\theta = 500$ W y olvida multiplicar por la velocidad.
    \item \textbf{C) 3\,460 W:} Usa $\cos 30^\circ \approx 0.866$ en lugar de $\cos 60^\circ = 0.5$ --- confunde el ángulo con su complementario.
    \item \textbf{D) 4\,000 W:} Calcula $P = F \times v = 1\,000 \times 4$ ignorando el ángulo.
\end{itemize}
\end{respuestabox}

% --- P7 ---
\begin{respuestabox}
\subsection{Pregunta 7}
\textcolor{verdecorrecto}{\textbf{Respuesta correcta: B) 10 m/s}}

\textbf{Justificación:}\\
\textbf{Fórmula utilizada:} Teorema Trabajo--Energía (Lección: Sección 3 --- Energía Cinética)
$$W = \frac{1}{2}mv_f^2 \quad \Rightarrow \quad v_f = \sqrt{\frac{2W}{m}} = \sqrt{\frac{2 \times 100\,\text{J}}{2\,\text{kg}}} = \sqrt{100\,\text{m}^2/\text{s}^2} = 10\,\text{m/s}$$

\textbf{Respuestas incorrectas:}
\begin{itemize}
    \item \textbf{A) 5 m/s:} Usa $v = \sqrt{W/m}$ sin el factor 2: $\sqrt{50} \approx 7$ m/s.
    \item \textbf{C) 50 m/s:} Calcula $v = W/m = 50$ --- olvida la raíz y el factor 2.
    \item \textbf{D) 100 m/s:} Toma el valor del trabajo directamente como velocidad.
\end{itemize}
\end{respuestabox}

% --- P8 ---
\begin{respuestabox}
\subsection{Pregunta 8}
\textcolor{verdecorrecto}{\textbf{Respuesta correcta: C) 20 m/s}}

\textbf{Justificación:}\\
\textbf{Fórmula utilizada:} Conservación de la energía mecánica (Lección: Sección 5)
$$v_f = \sqrt{2gh} = \sqrt{2(10)(20)} = \sqrt{400\,\text{m}^2/\text{s}^2} = 20\,\text{m/s}$$

La masa del clavadista se cancela; la velocidad no depende de ella.

\textbf{Respuestas incorrectas:}
\begin{itemize}
    \item \textbf{A) 10 m/s:} Usa $v = \sqrt{gh}$ sin el factor 2: $\sqrt{200} \approx 14$, redondeado.
    \item \textbf{B) 14 m/s:} Calcula $\sqrt{2gh}$ correctamente pero con $h = 10$ m en lugar de 20 m.
    \item \textbf{D) 200 m/s:} Calcula $2gh = 400$ sin aplicar la raíz cuadrada.
\end{itemize}
\end{respuestabox}

% --- P9 ---
\begin{respuestabox}
\subsection{Pregunta 9}
\textcolor{verdecorrecto}{\textbf{Respuesta correcta: B) 10 m/s}}

\textbf{Justificación:}\\
\textbf{Fórmulas:} $E_e = E_c$ (Lección: Secciones 4 y 5)

\textbf{Paso 1 --- Energía elástica:}
$$E_e = \frac{1}{2}(800\,\text{N/m})(0.5\,\text{m})^2 = \frac{1}{2}(800)(0.25) = 100\,\text{J}$$

\textbf{Paso 2 --- Velocidad de salida:}
$$v = \sqrt{\frac{2E_e}{m}} = \sqrt{\frac{2(100\,\text{J})}{2\,\text{kg}}} = \sqrt{100\,\text{m}^2/\text{s}^2} = 10\,\text{m/s}$$

\textbf{Respuestas incorrectas:}
\begin{itemize}
    \item \textbf{A) 5 m/s:} Olvida el factor 2: usa $\sqrt{E_e/m} = \sqrt{50} \approx 7$ m/s.
    \item \textbf{C) 20 m/s:} Calcula $2E_e/m = 100$ sin aplicar la raíz cuadrada.
    \item \textbf{D) 100 J:} Reporta la energía elástica como si fuera velocidad --- confunde unidades.
\end{itemize}
\end{respuestabox}

% --- P10 ---
\begin{respuestabox}
\subsection{Pregunta 10}
\textcolor{verdecorrecto}{\textbf{Respuesta correcta: B) 2 m/s}}

\textbf{Justificación:}\\
\textbf{Fórmula utilizada:} Conservación del ímpetu (Lección: Sección 6 --- Conservación del Ímpetu)
$$v_f = \frac{m_A v_{A,i}}{m_A + m_B} = \frac{(10\,000\,\text{kg})(4\,\text{m/s})}{20\,000\,\text{kg}} = 2\,\text{m/s}$$

\textbf{Respuestas incorrectas:}
\begin{itemize}
    \item \textbf{A) 4 m/s:} Conserva la velocidad original sin considerar el aumento de masa.
    \item \textbf{C) 1 m/s:} Divide la velocidad entre 4 en lugar de entre 2.
    \item \textbf{D) 0.5 m/s:} Divide entre 8 sin fundamento.
\end{itemize}
\end{respuestabox}

% --- P11 ---
\begin{respuestabox}
\subsection{Pregunta 11}
\textcolor{verdecorrecto}{\textbf{Respuesta correcta: B) 20 m/s}}

\textbf{Justificación:}\\
\textbf{Fórmula utilizada:} Conservación de la energía mecánica (Lección: Sección 5)
$$v_f = \sqrt{2gh} = \sqrt{2(10)(20)} = \sqrt{400\,\text{m}^2/\text{s}^2} = 20\,\text{m/s}$$

La masa del snowboarder no interviene.

\textbf{Respuestas incorrectas:}
\begin{itemize}
    \item \textbf{A) 10 m/s:} Usa $v = \sqrt{gh}$ sin el factor 2: $\sqrt{200} \approx 14$, redondeado.
    \item \textbf{C) 40 m/s:} Duplica el resultado correcto sin justificación.
    \item \textbf{D) 200 m/s:} Calcula $2gh = 400$ sin aplicar la raíz cuadrada.
\end{itemize}
\end{respuestabox}

% --- P12 ---
\begin{respuestabox}
\subsection{Pregunta 12}
\textcolor{verdecorrecto}{\textbf{Respuesta correcta: C) 65\%}}

\textbf{Justificación:}\\
\textbf{Fórmula utilizada:} Eficiencia (Lección: Sección 7 --- Procesos Disipativos)
$$\eta = \frac{W_{\text{útil}}}{W_{\text{entrada}}} \times 100\% = \frac{130\,\text{kJ}}{200\,\text{kJ}} \times 100\% = 65\%$$

\textbf{Respuestas incorrectas:}
\begin{itemize}
    \item \textbf{A) 35\%:} Calcula la fracción de energía \textit{perdida}: $(200-130)/200 = 35\%$ --- confunde pérdida con eficiencia.
    \item \textbf{B) 54\%:} Error aritmético al dividir $130/200$.
    \item \textbf{D) 154\%:} Invierte la fracción: $200/130$ --- divide entrada entre útil.
\end{itemize}
\end{respuestabox}

% --- P13 ---
\begin{respuestabox}
\subsection{Pregunta 13}
\textcolor{verdecorrecto}{\textbf{Respuesta correcta: B) 10 m/s}}

\textbf{Justificación:}\\
\textbf{Fórmulas:} Energía con fricción (Lección: Sección 7 --- Procesos Disipativos)

\textbf{Paso 1 --- Energía potencial inicial:}
$$E_{m,i} = mgh = 10\,\text{kg} \times 10\,\text{m/s}^2 \times 8\,\text{m} = 800\,\text{J}$$

\textbf{Paso 2 --- Calor disipado por fricción:}
$$Q = f \cdot d = 30\,\text{N} \times 10\,\text{m} = 300\,\text{J}$$

\textbf{Paso 3 --- Energía cinética final:}
$$E_{c,f} = 800\,\text{J} - 300\,\text{J} = 500\,\text{J}$$

\textbf{Paso 4 --- Velocidad final:}
$$v_f = \sqrt{\frac{2 \times 500\,\text{J}}{10\,\text{kg}}} = \sqrt{100\,\text{m}^2/\text{s}^2} = 10\,\text{m/s}$$

\textbf{Respuestas incorrectas:}
\begin{itemize}
    \item \textbf{A) 5 m/s:} Olvida el factor 2 al despejar $v$: usa $\sqrt{E_{c,f}/m} = \sqrt{50} \approx 7$ m/s.
    \item \textbf{C) 16 m/s:} No descuenta la fricción; usa $E_{c,f} = 800$ J directamente.
    \item \textbf{D) 20 m/s:} No descuenta fricción y además olvida el factor 2: $\sqrt{800/10} \approx 9$ --- o redondea erróneamente.
\end{itemize}
\end{respuestabox}

% --- P14 ---
\begin{respuestabox}
\subsection{Pregunta 14}
\textcolor{verdecorrecto}{\textbf{Respuesta correcta: C) 180 W}}

\textbf{Justificación:}\\
\textbf{Fórmulas:} Trabajo y potencia (Lección: Secciones 1 y 2)

\textbf{Paso 1 --- Trabajo contra la gravedad:}
$$W = mgh = 90\,\text{kg} \times 10\,\text{m/s}^2 \times 4\,\text{m} = 3\,600\,\text{J}$$

\textbf{Paso 2 --- Potencia:}
$$P = \frac{W}{t} = \frac{3\,600\,\text{J}}{20\,\text{s}} = 180\,\text{W}$$

\textbf{Respuestas incorrectas:}
\begin{itemize}
    \item \textbf{A) 18 W:} Olvida el factor $g$: calcula $mh/t = 90 \times 4 / 20 = 18$.
    \item \textbf{B) 72 W:} Error aritmético: divide $3\,600$ entre 50 en lugar de 20.
    \item \textbf{D) 3\,600 W:} Reporta el trabajo total como potencia --- olvida dividir entre el tiempo.
\end{itemize}
\end{respuestabox}

% --- P15 ---
\begin{respuestabox}
\subsection{Pregunta 15}
\textcolor{verdecorrecto}{\textbf{Respuesta correcta: C) 0.6 m/s}}

\textbf{Justificación:}\\
\textbf{Fórmula utilizada:} Conservación del ímpetu (Lección: Sección 6 --- Conservación del Ímpetu)
$$0 = m_{\text{astr.}} v_{\text{astr.}} + m_{\text{panel}} v_{\text{panel}}$$

\textbf{Despeje:}
$$v_{\text{astr.}} = \frac{-m_{\text{panel}} \cdot v_{\text{panel}}}{m_{\text{astr.}}} = \frac{-(20\,\text{kg})(3\,\text{m/s})}{100\,\text{kg}} = -0.6\,\text{m/s}$$

El signo negativo indica dirección opuesta al panel. Rapidez: 0.6 m/s.

\textbf{Verificación:} $p_{\text{total}} = (100)(-0.6) + (20)(3) = -60 + 60 = 0$ $\checkmark$

\textbf{Respuestas incorrectas:}
\begin{itemize}
    \item \textbf{A) 3 m/s:} Asume que ambos salen a la misma velocidad --- ignora la diferencia de masas.
    \item \textbf{B) 1.5 m/s:} Divide la velocidad del panel entre 2 sin justificación física.
    \item \textbf{D) 15 m/s:} Invierte la razón de masas: $v = (100)(3)/20 = 15$ m/s.
\end{itemize}
\end{respuestabox}
